% Root, B, and D literal audits passed; no new numerical gate.
\subsection{All five-state actions opposite a right-zero corner}

\begin{lemma}[The two eigenvalues of a component-preserving graph map]
\label{C:graph-component-eigenvalues}
Let a finite monoid $M$ act cellularly on a finite graph $\Gamma$,
put $T=\pi_0\Gamma$, and form the actual comparison
\[
 q:X=\operatorname{hocofib}(\Gamma_{hM}\to BM)
       \longrightarrow
 X_0=\operatorname{hocofib}(T_{hM}\to BM).
\]
Write $W=H_1(\Gamma;\mathbb Q)$ as a right $\mathbb QM$-module.
If $W\otimes_{\mathbb QM}\mathbb Q=0$, the pullback
$H^2(X_0;\mathbb Q)\to H^2(X;\mathbb Q)$ is surjective.

More generally, suppose an actual $M$-equivariant graph self-map
induces the identity on $T$ and multiplication by
$d\in\mathbb Z\setminus\{-1,0,1\}$ on $W$.
On every rational cohomology group of $X$, its induced map $f^*$
satisfies
\[
 (f^*-1)(f^*-d)=0,
 \qquad
 \ker(f^*-1)=\operatorname{im}\bigl(q^*:H^*(X_0;\mathbb Q)
                                   \to H^*(X;\mathbb Q)\bigr).
\]
If $\dim H^2(X;\mathbb Q)=1$ and all degree-two products on $X_0$
vanish, then the degree-two square on $X$ vanishes as well.
\end{lemma}
\begin{proof}
Let $A_\Gamma$ be the augmented graph complex with
$C_1(\Gamma;\mathbb Q)$ in degree one and
$\widetilde C_0(\Gamma;\mathbb Q)$ in degree zero.
Projection to the component set gives a degreewise surjection
of right-module complexes
\[
 A_\Gamma\longrightarrow\widetilde{\mathbb QT}.
\]
Its kernel has $C_1(\Gamma;\mathbb Q)$ in degree one and
$\ker(C_0(\Gamma;\mathbb Q)\to\mathbb QT)$ in degree zero.
This kernel has only the homology $W$ in degree one; the literal
inclusion of graph cycles is a natural equivariant quasi-isomorphism.

Tensor this short exact sequence with the free left monoid bar
resolution and suspend once. These are natural reduced chain models
for the two cofibers and their actual map $q$. The kernel is
quasi-isomorphic to the twice-suspended two-sided bar of $W$.
Consequently the cohomology exact sequence has the segment
\[
 H^2(X_0;\mathbb Q)\longrightarrow H^2(X;\mathbb Q)
 \longrightarrow
 \operatorname{Hom}_{\mathbb Q}
       (W\otimes_{\mathbb QM}\mathbb Q,\mathbb Q),
\]
which proves the first assertion.

For the self-map, the action on the component complex is the
identity and its action on the kernel's bar homology is scalar $d$.
This is a natural statement on the actual chain models: graph
cycles have no two-dimensional boundaries, so their induced map
is literally multiplication by $d$. In the cohomology exact
sequence, $f^*-d$ therefore carries every class on $X$ into the
image of $q^*$, where $f^*$ is the identity. This proves the
polynomial identity. A class fixed by $f^*$ maps to a class both
fixed and multiplied by $d$ in the kernel term; since $d\ne1$,
that image is zero. Exactness gives the stated eigenvalue-one
pullback assertion.

On the one-dimensional space $H^2(X;\mathbb Q)$, $f^*$ is scalar
either $1$ or $d$. In the first case its generator pulls back from
$X_0$, so its square is zero by hypothesis. In the second case its
square is multiplied by $d^2$. The displayed polynomial in degree
four annihilates this square, and
$(d^2-1)(d^2-d)\ne0$. The square is again zero.
This argument makes no finite-type inference about $X_0$.
\end{proof}

\begin{theorem}[Five-state right-zero reduction]
\label{C:right-zero-five-state-reduction}
Let $R=\{1,p,q\}$, with $ab=b$ for $a,b\in\{p,q\}$, and let
$S$ be an $(R,M)$-biset with $|S|\leq5$.
If its two-object collage has classifying space $\mathbb CP^2$,
then an ordinary right $M$-act cone on at most $|S|$ states has
that homotopy type.
\end{theorem}
\begin{proof}
Use the two retractions $P,Q:S\to U$ and their natural right
$M$-action graph $\Gamma$. Its collage is
$X=\operatorname{hocofib}(\Gamma_{hM}\to BM)$.
The connected case is excluded by
Theorem~\ref{C:right-zero-five-connected}.

If $\Gamma$ is a forest, its projection to the component set
$T$ is an equivariant underlying homotopy equivalence. The bar
construction preserves it, so $X$ is actually equivalent to the
ordinary $T$-act cone. Since $|T|\leq|S|$, this is the required
reduction.

If $P=Q$, their common retraction is an actual biset retraction
$S\to U$. The collage $Y$ on $U$ is a homotopy retract of $X$
and is equivalent to the ordinary $U$-act cone, because the left
$R$-action there is trivial and $BR$ is contractible.
The relative quotient of the two double-mapping-cylinder models
is $\Sigma(E_S/E_U)$, precisely as in the $P=Q$ part of
Theorem~\ref{C:right-zero-four-state-reduction}.
If $H^2(Y;\mathbb Q)=0$, all degree-two classes of $X$ are relative
and their squares vanish. Otherwise the split integral homology
inclusion into $H_2(X;\mathbb Z)=\mathbb Z$ is an isomorphism.
The nonzero square forces the same conclusion in degree four.
Every other reduced homology group of $Y$ is zero because $Y$
is a retract of $X\simeq\mathbb CP^2$, and its fundamental group
is trivial for the same reason. The inclusion $Y\to X$ is an
integral homology equivalence between simply connected CW spaces,
so it is a homotopy equivalence. This obtains the actual ordinary
act of the assertion; it does not replace finite type by rational
homology alone.

It remains that $\Gamma$ is disconnected, is not a forest, and
has at least one nonloop edge, since $P\ne Q$.
The bound $|U|+|S\setminus U|\leq5$ forces exactly three
vertices and two edges. There is one cyclic component and one
acyclic component, so $T$ has two points and $W=H_1(\Gamma;\mathbb Q)$
has dimension one. The cyclic component is either a two-edge circle,
a loop with a connecting tail, or an isolated loop with a separate
interval component. These exhaust two edges in a disconnected graph
with a nonloop edge and a cycle.

If some actor has degree different from $+1$ on $W$, its scalar
relation kills $W\otimes_{\mathbb QM}\mathbb Q$.
Lemma~\ref{C:graph-component-eigenvalues} then makes every degree-two
class pull back from the ordinary two-state $T$-act cone, whose
products vanish by Lemma~\ref{C:finite-line-act-cups}.
We may therefore assume that every actor induces degree $+1$.
Every actor must preserve the unique cyclic component; mapping it
to the other component would induce degree zero.

When that component contains a loop edge, degree $+1$ forces
the unique loop edge to map to itself, hence forces its vertex
to be fixed by all of $M$. That vertex is also fixed by $R$.
The opposite form of Lemma~\ref{C:fixed-incoming-cups} makes
every degree-two product zero, a contradiction.

The only remaining case is a two-edge circle and an isolated
vertex. If the circle edges have the same orientation, a degree-one
actor fixes both vertices and both edges; an edge interchange has
degree minus one and a fold has degree zero. If they have opposite
orientations, a degree-one actor is either the identity or the
half-rotation interchanging both vertices and edges.
Identify the circle with the unit circle so its two vertices are
$1$ and $-1$. The map $z\mapsto z^3$ fixes these vertices and
commutes with every permitted actor on the cyclic component.
Extend it by the identity on the isolated vertex. An actor can
send that isolated vertex only to a vertex of the graph, all of
which the self-map fixes. Thus this is a strictly $M$-equivariant
graph map, inducing the identity on $T$ and multiplication by
three on $W$.

The ordinary two-state $T$-act cone has zero degree-two products.
Lemma~\ref{C:graph-component-eigenvalues}, with $d=3$ and
$\dim H^2(X;\mathbb Q)=1$, contradicts the projective-plane square.
Thus the remaining disconnected cases cannot occur, leaving only
the actual ordinary-act reductions above.
\end{proof}

\begin{corollary}[The exact small-state right-zero-corner minimum]
\label{C:right-zero-five-state-minimum24}
Among triangular two-object categories with the fixed corner
$R=\{1,p,q\}$ above and at most five cross arrows, the least
arrow count realizing $\mathbb CP^2$ is twenty-four.
\end{corollary}
\begin{proof}
The preceding theorem gives an ordinary act cone with the same
source monoid and at most as many states. Its arrow count is at
least twenty-two by Theorem~\ref{C:act-cone-minimum-twenty-two}.
Restoring the three-element corner in place of the singleton adds
two arrows, so the original category has at least twenty-four.

For equality take either twenty-two-arrow act-cone model of
Theorem~\ref{C:CP2-upper22}, with an eighteen-element source
monoid and three states. Replace its trivial sink corner by $R$
and give all elements of $R$ the identity action on those states.
Its edge action category is the disjoint union of copies of $R$;
the projection to the state set is equivariantly a homotopy
equivalence because $BR$ is contractible. The collage is therefore
equivalent to the original act cone, and has $18+3+3=24$ arrows.
\end{proof}
